%%%%%%%%%%%%%%%%%%%%%%%%%%%%%%%%%%%%%%%%%
% Cv BENGO
% LuaLaTeX Template
% Version 1.12.0 (24/04/2023)
%
% Original authors:
% Ged Lex (gedlex@hotmail.ch)
%
% Important note:
% This template must be compiled with LuaLaTeX, the below lines will ensure this
%!TEX TS-program = lualatex
%!TEX encoding = UTF-8 Unicode
%%%%%%%%%%%%%%%%%%%%%%%%%%%%%%%%%%%%%%%%%

%----------------------------------------------------------------------------------------
%	PACKAGES AND OTHER DOCUMENT CONFIGURATIONS

%----------------------------------------------------------------------------------------
\documentclass[rightPos]{ReCeiVe}      % By default, use 'letterpaper' for US letter
%\geometry{letterpaper, % paper size
%          left=1.75cm, right=1.75cm, top=1.5cm, bottom=1.5cm, footskip=.6cm, headsep=.5cm, % margins
%          showframe}                  % show geometry frame
%          heightrounded}              % avoid underfull vbox warning
 \usepackage{xcolor} 
 \usepackage{multicol}
 \usepackage{comment}
% Color for highlights
%\definecolor{highlight}{HTML}{FF0000} % Specify your own color
%\colorlet{highlight}{white}           % Set predefined color
% Default colors include: darkgray, gray, lightgray, lightblue, orange, red, concrete

%Colors for text - uncomment and modify
%\definecolor{darktext}{HTML}{FF0000}
%\definecolor{text}{HTML}{FF0000}
%\definecolor{graytext}{HTML}{FF0000}
%\definecolor{lighttext}{HTML}{414141}

%----------------------------------------------------------------------------------------
%	PERSONAL INFORMATION
%	Comment any of the lines below if they are not required
%----------------------------------------------------------------------------------------
\background{pics/background5.pdf}
\photo[rectangle,edge,nofill,left]{5cm}{pics/pp2.png}
\name{Dylane Bengono}{}
%\address{}
\homepage{dylane-mu.vercel.app/}
\mobile{+237 691457703}
\email{chaneldylanebengono@gmail.com}
\linkedin{chanel-dylane-b-91b850194/}

\extrainfo{\textit{Permis de conduire} \textbf{\textit{ Cat A  \& B}}} 

%\twitter{@twit}
%\xing{xing name}
%\stackoverflow{SOid}{SOname}
%\skype{skypeid}
%\reddit{reddit account}
%\extrainfo{info}

\position{Responsable Sécurité des Systèmes d'Information (RSSI) | Expert Cybersécurité}
\headwords{Gouvernance SI | ISO 27001 | Pentest | SOC/SIEM | DevSecOps | Cloud Security | EBIOS RM | PCA/PRA}

%\quote{"Certains voient les choses telles qu'elles sont et se demandent pourquoi. Moi, je rêve de choses qui n'ont jamais existé et je me demande pourquoi pas." - Robert Kennedy}

%----------------------------------------------------------------------------------------
%	SIDEBAR CONTENT
%	Fill in the information you would like to add to your sidebar
%----------------------------------------------------------------------------------------

\aboutMe{
RSSI à la BVMAC, Ingénieur Polytechnicien (BAC+5) spécialisé en Cybersécurité \& Investigation Numérique.
Architecte de la sécurité des SI : conception de PSSI, déploiement de SOC/SIEM, pilotage PCA/PRA et gestion des risques (ISO 27001/27005, EBIOS RM, RGPD).
Expert en sécurité offensive (pentest Black/Grey/White box, Red Team) et défensive (SOC, CSIRT, Threat Hunting, Forensics).
Maîtrise des environnements Cloud (AWS, GCP, Azure) et des architectures DevSecOps (Kubernetes, Terraform, CI/CD).
Certifié ISO 27001, CCNP, NSE4 Fortinet, CPTA v2, CCT EC-Council.
Lauréat national en cybersécurité \& Hall of Fame EC-Council.
}

\skillset[label/0]{
Gouvernance \& PSSI/.95,
Gestion des Risques/.9,
Conformité ISO 27001/.9,
SOC \& SIEM/.85,
Pentest \& Red Team/.85,
Sécurité Réseau/.9,
Cloud \& DevSecOps/.8
}

\languages{{Anglais/C1}, {Français/C2}}

%\hobbies{Jeux vidéo, Voyage, Volleyball}


% Usage: \userSection[<section title>]{<content>}
% \userSection[title]{content}

%----------------------------------------------------------------------------------------
\begin{document}
% Print the header
% Usage: \makecvheader[<position>]
\makecvheader
% Print the sidebar
% Usage: \makecvsidebar[xOffset/value,yOffset/value,noRadius|radius]
\makecvsidebar
%----------------------------------------------------------------------------------------
%	CV/RESUME CONTENT
%	Each section is imported separately, open each file in order to modify it
%----------------------------------------------------------------------------------------
\section{Expérience Professionnelle}


\cventry
{Responsable Sécurité des Systèmes d'Information \textbf{| RSSI}} % Titre
{Bourse des Valeurs Mobilières de l'Afrique Centrale \textbf{| BVMAC}} % Organisation
{CDI} % Contrat
{Depuis Oct. 2024} % Date
\begin{cvitems}
\item {\textbf{Gouvernance SI :} Élaboration et déploiement de la PSSI, cartographie complète des actifs et classification des données sensibles.}
\item {\textbf{Résilience :} Conception et mise en œuvre du PCA et du PRA ; pilotage des tests de reprise pour garantir la continuité opérationnelle.}
\item {\textbf{SOC/SIEM :} Déploiement et administration d'une solution SIEM (Wazuh), définition des cas d'usage, corrélation des alertes et réponse aux incidents 24/7.}
\item {\textbf{Pentest :} Planification et réalisation de tests d'intrusion (Black/Grey/White Box) sur le SI, les applicatifs et les réseaux (BVMAC + sites de réplication).}
\item {\textbf{Supervision réseau :} Mise en place d'outils de monitoring centralisés couvrant le réseau BVMAC, les SDBS et le site de repli.}
\item {\textbf{Conformité :} Pilotage des actions de conformité RGPD, recommandations COBAC et guide d'hygiène ANSSI ; gestion des risques (ISO 27005, EBIOS RM).}
\item {\textbf{BYOD \& IAM :} Implémentation d'une politique BYOD et d'un cadre de gestion des accès (moindre privilège, MFA, révocation des droits).}
\item {\textbf{Sauvegarde :} Définition et déploiement d'une stratégie de sauvegarde et de réplication des données critiques de l'infrastructure financière.}
\end{cvitems}


\cventry
{Security Engineer \textbf{| Ingénieur Sécurité} — Full Remote} % Titre
{adorsys GmbH \& Co. KG \textbf{| Fintech, Allemagne}} % Organisation
{CDI} % Contrat
{Depuis Déc. 2023} % Date
\begin{cvitems}
\item {\textbf{Open Source :} Développement d'une bibliothèque Java pour SD-JWT (identité numérique décentralisée) — contribution au projet adorsys/sd-jwt sur GitHub.}
\item {\textbf{IDS/IPS :} Automatisation du déploiement de Snort 3 avec Wazuh sur Kubernetes pour la détection et la prévention d'intrusions réseau (NFV).}
\item {\textbf{EPP/XDR/SIEM :} Déploiement d'une plateforme unifiée de protection des endpoints et de corrélation des événements de sécurité.}
\item {\textbf{SMSI :} Mise en place et maintien d'un Système de Management de la Sécurité de l'Information (SMSI) conforme à l'ISO 27001.}
\item {\textbf{DevSecOps :} Industrialisation des déploiements (Helm, Terraform, Tekton CI/CD, ArgoCD) sur AWS et GCP — intégration de la sécurité dès la conception.}
\end{cvitems}

\cventry
{Ingénieur Cybersécurité \textbf{| Projet de Fin d'Études}} % Titre
{Port Autonome de Kribi \textbf{| PAK}} % Organisation
{PFE} % Contrat
{Fév. 2023 — Nov. 2023} % Date
\begin{cvitems}
\item {\textbf{SIEM :} Analyse, conception et déploiement de Wazuh pour la supervision et la corrélation des événements de sécurité sur l'infrastructure portuaire.}
\item {\textbf{Détection :} Développement de règles personnalisées et de décodeurs adaptés aux flux spécifiques du SI (logs réseau, systèmes, applicatifs).}
\item {\textbf{IDS/IPS :} Intégration de Snort pour la détection des menaces sur le périmètre réseau.}
\end{cvitems}

\cventry
{Ingénieur Sécurité Réseaux \textbf{| Stage}} % Titre
{Délégation Générale à la Sûreté Nationale \textbf{| DGSN}} % Organisation
{Stage} % Contrat
{2022} % Date
\begin{cvitems}
\item {\textbf{NAC :} Conception et déploiement d'un contrôle d'accès réseau (Network Access Control) pour sécuriser les accès au SI.}
\item {\textbf{Audit :} Cartographie des actifs SI, analyse des risques et mise en place d'une politique de surveillance du réseau.}
\end{cvitems}

\vspace{3mm}
% ----- À PARTIR D'ICI ON SUPPRIME LA SIDEBAR -----
\clearpage
\newgeometry{left=2cm, right=2cm, top=2cm, bottom=2cm}
% Supprimer l'espace avant la première section
\vspace{-\acvSectionTopSkip}
\section{Formation}
\cventry
{Master of Engineering — Cybersécurité \& Investigation Numérique} % Titre
{École Nationale Supérieure Polytechnique de Yaoundé | ENSPY} % Établissement
{Université de Yaoundé I} % Lieu
{BAC+5 | 2023} % Date
\begin{cvitems}
\item {Spécialisation : Sécurité des réseaux, Cryptographie, Forensics, Gouvernance SI, Gestion des risques, Droit du numérique.}
\end{cvitems}

\vspace{3mm}

\section{Compétences Techniques}
{\color{blue} \rule{\linewidth}{1mm}}
\vspace{-3mm}
\begin{multicols}{2}

\subsection*{Gouvernance \& Conformité}
\begin{itemize}
  \item ISO 27001, 27002, 27005 | EBIOS RM
  \item RGPD | Guide ANSSI | NIST CSF
  \item PSSI, SMSI, PCA / PRA
  \item Audit SI, Analyse de risques
\end{itemize}

\subsection*{Sécurité Offensive}
\begin{itemize}
  \item Pentest (Black / Grey / White Box)
  \item Kali Linux, Metasploit, Nmap, Nessus, Burp Suite
  \item OSINT, Ingénierie sociale
  \item CTF — Forensics, Reverse, Web, Crypto
\end{itemize}

\columnbreak

\subsection*{Sécurité Défensive \& SOC/SIEM}
\begin{itemize}
  \item Wazuh | Splunk | IBM QRadar | Snort 3
  \item EPP / EDR / XDR | IDS / IPS
  \item Threat Hunting, Incident Response, Malware Triage
  \item Digital Forensics \& Investigation Numérique
\end{itemize}

\subsection*{Infrastructure, Cloud \& DevSecOps}
\begin{itemize}
  \item Kubernetes, Docker, Helm, Terraform, Ansible
  \item Tekton CI/CD | ArgoCD | GitHub Actions
  \item AWS, GCP, Azure
  \item IAM, Zero Trust, PacketFence NAC
\end{itemize}

\subsection*{Réseaux \& Systèmes}
\begin{itemize}
  \item Firewall, VPN, MPLS, TCP/IP, DNS, DHCP
  \item Cisco (CCNP) | Fortinet (NSE4) | Netbox
  \item Linux (Admin avancée), Windows Server
  \item Supervision : Zabbix, Nagios
\end{itemize}

\subsection*{Développement \& Scripting}
\begin{itemize}
  \item Python, Bash, Go, Rust
  \item Java, JavaScript / Node.js, React
  \item Cryptographie : RSA, AES, SHA, TLS/SSL, IPsec
\end{itemize}

\end{multicols}

\vspace{3mm}
\section{Projets Techniques \& Open Source}
{\color{blue} \rule{\linewidth}{1mm}}
\vspace{-3mm}
\begin{multicols}{2}

\textbf{Sécurité Réseau}
\begin{itemize}
  \item \textbf{snort-docker} — Conteneurisation de Snort 3 pour le déploiement en environnement NFV ; images Docker publiées sur GitHub.
  \item \textbf{snort3\_aws} — Déploiement IDS/IPS Snort 3 sur AWS (GWLB + EKS) avec orchestration Kubernetes ; architecture scalable haute disponibilité.
  \item \textbf{Wazuh Kubernetes} — Déploiement de la plateforme XDR/SIEM Wazuh sur cluster K8s ; intégration avec Snort pour une défense en profondeur.
\end{itemize}

\columnbreak

\textbf{Identité Numérique \& Conformité}
\begin{itemize}
  \item \textbf{adorsys/sd-jwt} — Contribution à la bibliothèque Java pour la production et vérification de SD-JWT (standard IETF pour l'identité décentralisée).
  \item \textbf{keycloak-ssi-deployment} — Configuration de déploiement Keycloak pour la gestion des identités et accès (SSI/SSO).
  \item \textbf{GRC Framework} — Développement d'un cadre de conformité supportant 70+ référentiels de sécurité (ISO, NIST, RGPD...).
\end{itemize}

\end{multicols}

\vspace{3mm}
\section{Distinctions \& Récompenses}
\begin{cvhonors}
\cvhonor{LAURÉAT}{Programme d'Excellence o'100 THE OKWELEANS 3\textsuperscript{e} Édition}{Yaoundé, Cameroun}{2023}
\cvhonor{Hall of Fame}{EC-Council Cyber Challenge — Reconnaissance internationale}{EC-COUNCIL}{2022}
\cvhonor{LAURÉAT}{Olympiades des Experts en Cybersécurité | AEC-CTF}{En Ligne}{2021}
\end{cvhonors}

\section{Certifications}
{\color{blue} \rule{\linewidth}{1mm}}
\vspace{-3mm}
\begin{multicols}{2}
\textbf{Gouvernance \& Conformité}
\begin{itemize}
  \item ISO/IEC 27001 Associate — PECB (2024)
  \item Information Systems Auditing, Controls and Assurance
  \item MOOC Sécurité des Systèmes — ANSSI
\end{itemize}

\medskip
\textbf{Réseaux \& Infrastructure}
\begin{itemize}
  \item CCNP — Cisco Certified Network Professional (2024)
  \item Fortinet NSE4 — Network Security Professional (2022)
  \item Cisco Ethical Hacker
\end{itemize}

\columnbreak

\textbf{Cybersécurité Offensive \& Défensive}
\begin{itemize}
  \item Certified API Security Analyst (CASA) — APIsec (2026)
  \item Certified Cybersecurity Technician (CCT) — EC-Council (2024)
  \item Certified Purple Team Analyst v2 (CPTA) — (2023)
  \item Penetration Testing, Incident Response and Forensics — IBM
  \item Certified Secure Computer User V2 (CSCU)
  \item Cybersecurity and the Internet of Things
\end{itemize}

\medskip
\textbf{Cloud \& DevOps}
\begin{itemize}
  \item Azure DevOps
  \item Python Developer
  \item Validation des Compétences en TIC
\end{itemize}

\end{multicols}

% ----- RESTAURATION DE LA GÉOMÉTRIE INITIALE -----
\restoregeometry


%----------------------------------------------------------------------------------------
\end{document}

